\chapter{Аналитический раздел}

В разделе формализуются объекты синтезируемой сцены, рассматриваются методы решения задач, перечисленных во введении, и выбираются алгоритмы для реализации.

\section{Формализация объектов сцены}

Линейные размеры в сцене задаются в условных единицах (у.е.). Ось $z$ направлена вертикально вверх, плоскость $xOy$ горизонтальна. Сцена занимает по горизонтали квадрат $[0;\,100] \times [0;\,100]$~у.е.

Сцена содержит следующие объекты.
\begin{enumerate}
	\item Ландшафт "--- заснеженная горная поверхность, заданная картой высот. Карта высот "--- полутоновое изображение, яркость пикселя которого линейно переводится в высоту поверхности $H(x, y) \in [0;\,10]$~у.е. Карта высот загружается из файла описания сцены и пользователем не изменяется.
	\item Объекты, размещаемые на ландшафте, количество которых $M \geqslant 5$. Каждый объект относится к одному из заранее подготовленных типов: домик, ель, автомобиль. Тип объекта составлен из стандартных тел (параллелепипеда, пирамиды, цилиндра и конуса), каждое из которых аппроксимируется многогранником; состав типов задаётся файлом описания сцены. Объект характеризуется:
	\begin{itemize}
		\item типом;
		\item положением $(x_i, y_i)$ на горизонтальной плоскости;
		\item углом поворота $\psi_i$ вокруг вертикальной оси.
	\end{itemize}
	Высота основания объекта равна $H(x_i, y_i)$, то есть объект стоит на поверхности ландшафта. Пользователь может добавлять, перемещать и удалять объекты при условии $M \geqslant 5$. Начальная расстановка объектов задаётся файлом описания сцены.
	\item Источник света на бесконечности, описываемый:
	\begin{itemize}
		\item направлением, которое задаётся азимутом $\varphi_L$ и высотой над горизонтом $\theta_L$; единичный вектор направления на источник
		\begin{equation}
			\vec{L} = (\cos\theta_L \cos\varphi_L,\ \cos\theta_L \sin\varphi_L,\ \sin\theta_L);
		\end{equation}
		\item цветом $C_L$;
		\item интенсивностью $I_L$.
	\end{itemize}
	\item Занавесы северного сияния, количество которых $K \leqslant 8$. Основание занавеса "--- кривая на горизонтальной плоскости, отклоняющаяся от своего основного направления на величину порядка $A$ с характерной длиной складки $\lambda$. Занавес "--- вертикальный светящийся слой над основанием, ограниченный по высоте нижней $h_{\min}$ и верхней $h_{\max}$ границами свечения. Каждый занавес характеризуется:
	\begin{itemize}
		\item высотами границ свечения $h_{\min}$ и $h_{\max}$;
		\item протяжённостью $\ell$ "--- длиной занавеса вдоль основного направления;
		\item складчатостью, которая задаётся амплитудой складок $A$ и характерной длиной складки $\lambda$;
		\item цветами свечения $C_{\text{н}}$, $C_{\text{с}}$ и $C_{\text{в}}$ нижней, средней и верхней частей занавеса; части получаются делением отрезка $[h_{\min};\, h_{\max}]$ на три равные доли.
	\end{itemize}
	Положение и основное направление каждого занавеса порождаются генератором псевдослучайных чисел.
	\item Светящаяся среда, общая для всех занавесов, описываемая:
	\begin{itemize}
		\item коэффициентом излучения $\varepsilon$;
		\item коэффициентом поглощения $\kappa$;
		\item начальным значением генератора псевдослучайных чисел $r_0$.
	\end{itemize}
	\item Наблюдатель, описываемый положением $P_V$ и направлением взгляда, которое задаётся азимутом $\varphi_V$ и углом места $\theta_V$. Угол обзора постоянен.
	\item Расчётная сетка размером $N_x \times N_y \times N_z$ узлов, покрывающая область $[0;\,100] \times [0;\,100] \times [10;\,60]$~у.е., в которой могут находиться занавесы, и шаг интегрирования $\Delta s$ вдоль луча.
\end{enumerate}

Оптические свойства поверхностей ландшафта и каждого типа объектов описываются коэффициентом диффузного отражения $k_d$, коэффициентом зеркального отражения $k_s$ и показателем блеска $n$; цвета поверхностей задаются файлом описания сцены. Пользователь может отключить отображение ландшафта и объектов, оставив только северное сияние.

Все перечисленные параметры, кроме содержимого файла описания сцены, задаются пользователем во время работы программы. Допустимые диапазоны параметров приведены в таблице~\ref{tab:input-ranges}.

\begin{table}[H]
	\caption{Допустимые диапазоны входных параметров}
	\label{tab:input-ranges}
	\begin{tabular}{|L{0.21}|L{0.33}|C{0.19}|C{0.26}|}
		\hline
		Объект & Параметр & Обозначение & Диапазон \\
		\hline
		Объекты & количество & $M$ & $M \geqslant 5$ \\
		\cline{2-4}
		& положение & $x_i$, $y_i$ & $[0;\,100]$ \\
		\cline{2-4}
		& угол поворота & $\psi_i$ & $[0^\circ;\,360^\circ)$ \\
		\hline
		Поверхности & коэффициенты отражения & $k_d$, $k_s$ & $[0;\,1]$, $k_d + k_s \leqslant 1$ \\
		\cline{2-4}
		& показатель блеска & $n$ & $[1;\,100]$ \\
		\hline
		Источник света & азимут, высота & $\varphi_L$, $\theta_L$ & $[0^\circ;\,360^\circ)$, $[0^\circ;\,90^\circ]$ \\
		\cline{2-4}
		& цвет & $C_L$ & RGB, $[0;\,1]^3$ \\
		\cline{2-4}
		& интенсивность & $I_L$ & $[0;\,1]$ \\
		\hline
		Занавесы & количество & $K$ & $[1;\,8]$ \\
		\cline{2-4}
		& границы свечения & $h_{\min}$, $h_{\max}$ & $10 \leqslant h_{\min} < h_{\max} \leqslant 60$ \\
		\cline{2-4}
		& протяжённость & $\ell$ & $[5;\,100]$ \\
		\cline{2-4}
		& амплитуда складок & $A$ & $[0;\,20]$ \\
		\cline{2-4}
		& длина складки & $\lambda$ & $[1;\,\ell]$ \\
		\cline{2-4}
		& цвета частей & $C_{\text{н}}$, $C_{\text{с}}$, $C_{\text{в}}$ & RGB, $[0;\,1]^3$ \\
		\hline
		Светящаяся среда & коэффициенты излучения и поглощения & $\varepsilon$, $\kappa$ & $[0;\,10]$, $[0;\,1]$ \\
		\cline{2-4}
		& начальное значение генератора & $r_0$ & $[0;\,2^{32} - 1]$ \\
		\hline
		Наблюдатель & положение & $P_V$ & внутри сцены, $z_V > H(x_V, y_V)$ \\
		\cline{2-4}
		& азимут, угол места & $\varphi_V$, $\theta_V$ & $[0^\circ;\,360^\circ)$, $[-90^\circ;\,90^\circ]$ \\
		\hline
		Расчёт & размер сетки & $N_x$, $N_y$, $N_z$ & $[16;\,256]$, $[16;\,256]$, $[8;\,128]$ \\
		\cline{2-4}
		& шаг интегрирования & $\Delta s$ & $[0{,}05;\,5]$ \\
		\hline
	\end{tabular}
\end{table}

Ограничения $M \geqslant 5$, $K \leqslant 8$ и предельный размер сетки $256 \times 256 \times 128$ узлов установлены заданием. Остальные ограничения обоснованы следующим образом.
\begin{itemize}
	\item Нижняя граница свечения не ниже наибольшей высоты ландшафта, $h_{\min} \geqslant 10$, поэтому занавесы не пересекают рельеф.
	\item Сумма $k_d + k_s$ не превышает единицы, так как поверхность не может отразить больше света, чем на неё падает.
	\item Высота источника над горизонтом неотрицательна, так как источник под горизонтом не освещает сцену.
	\item Длина складки ограничена снизу разрешением расчётной сетки. По теореме отсчётов (теореме Котельникова) периодическая составляющая, дискретизованная с шагом $\delta$, передаётся без искажений, только если её период больше $2\delta$~\cite{gonzalez2012}. Шаг сетки по горизонтали $\delta = 100 / N_x$, на сетке наибольшего размера $\delta \approx 0{,}39$~у.е., откуда $\lambda > 0{,}78$~у.е.; поэтому принято $\lambda \geqslant 1$~у.е. На более грубой сетке складки короче $2\delta$ будут искажаться.
	\item Диапазон шага интегрирования охватывает значения от много меньших шага сетки наибольшего размера до одной десятой толщины слоя, в котором находятся занавесы.
\end{itemize}

Формализация поставленной задачи в нотации IDEF0 представлена на рисунке~\ref{fig:idef0-a0}.

\begin{figure}[H]
	\centering
	\IfFileExists{img/idef0-a0.pdf}{%
		\includegraphics[width=\linewidth]{img/idef0-a0.pdf}%
	}{%
		\fbox{\parbox[c][6cm][c]{0.9\linewidth}{\centering Диаграмма IDEF0 A-0 (экспорт из Ramus: img/idef0-a0.pdf)}}%
	}
	\caption{Функциональная модель программы в нотации IDEF0}
	\label{fig:idef0-a0}
\end{figure}
